% Pack cover (doc.py bundle): the front cover and the "In this pack" page of several documents in one PDF.
% A pack is a stack of documents, so the cover shows the stack: fanned sheets on the band, one per part with a
% strip in that part's type colour; under the band a strip split into the parts in proportion to their pages
% (the spine is split the same way); and the parts as linked cards pinned above the footer.
% bundle.py sets, before cover-kit.tex: \PackTypeColours (\definecolor{PackType<n>} for every part, print
% greys included), \PackSegments (n/from/to as fractions of the pack's pages), \PackSheets (1,...,k: the
% sheets drawn, at most five), \PackCards (\packcard{n}{title}{type}{first}{last}{pages}) and
% \PackMap (\packmappart ... \packmapcols, the parts with their outline). Every part links to its first page
% through the named destination bundle-part-<n>; a map entry through bundle-p<n>-<k>.
\providecommand{\PackTypeColours}{}
\providecommand{\PackSegments}{1/0/1}
\providecommand{\PackSheets}{1}
\providecommand{\PackCards}{}
\providecommand{\PackMap}{}
\PackTypeColours

% ---------- front cover ----------
\thispagestyle{empty}
\pgfmathtruncatemacro{\packk}{0}\foreach \s in \PackSheets {\xdef\packk{\s}}
\begin{tikzpicture}[remember picture, overlay]
	% faint pages in the white middle, in the parts' colours, fading out towards the title (like the other covers)
	\begin{scope}[shift={(current page.south west)}]
		\foreach \row/\y/\dx in {0/8.7/0.55, 1/10.7/0, 2/12.7/0.85} \foreach \x [count=\i] in {5.0,6.75,...,19.5} {
			\fadeop{\x+\dx}{0.42}
			\pgfmathtruncatemacro{\c}{mod(\i+\row,\packk)+1}
			\pgfmathsetmacro{\tilt}{12*sin(\i*67+\row*41)}
			\begin{scope}[shift={(\x+\dx,\y)}, rotate=\tilt, opacity=\op]
				\draw[PackType\c, line width=0.9pt, line join=round] (0,0) -- (1.1,0) -- (1.1,1.1) -- (0.75,1.45) -- (0,1.45) -- cycle;
				\draw[PackType\c, line width=0.7pt] (0.75,1.45) -- (0.75,1.1) -- (1.1,1.1);
				\foreach \l/\w in {0.95/0.5, 0.7/0.75, 0.45/0.6, 0.2/0.7} \draw[PackType\c, line width=0.9pt, line cap=round] (0.18,\l) -- ++(\w,0);
			\end{scope}}
	\end{scope}
	% the band
	\fill[CoverBand] (current page.north west) rectangle ($(current page.north east)+(0,-9.55cm)$);
	% the parts as fanned sheets on the band, back to front; each has a strip in its type colour
	\begin{scope}[shift={(current page.north west)}]
		\foreach \s in \PackSheets {
			\pgfmathsetmacro{\ang}{(\packk+1)/2*9-\s*9}
			\begin{scope}[rotate around={\ang:(17.4,-8.6)}]
				\fill[CoverBand] (15.55,-8.3) rectangle (19.25,-3.4);
				\fill[CoverMotifOnBand, opacity=0.10] (15.55,-8.3) rectangle (19.25,-3.4);
				\fill[PackType\s] (15.55,-3.95) rectangle (19.25,-3.4);
				\draw[CoverMotifOnBand, opacity=0.6, line width=0.9pt] (15.55,-8.3) rectangle (19.25,-3.4);
			\end{scope}}
		% lines of text on the front sheet
		\pgfmathsetmacro{\ang}{(\packk+1)/2*9-\packk*9}
		\begin{scope}[rotate around={\ang:(17.4,-8.6)}]
			\foreach \y/\w in {4.6/2.6, 5.1/3.0, 5.6/2.2, 6.1/2.8, 6.6/1.8, 7.1/2.5}
				\draw[CoverMotifOnBand, opacity=0.35, line width=1.6pt, line cap=round] (15.95,-\y) -- ++(\w,0);
		\end{scope}
	\end{scope}
	% under the band: the pack split into its parts, in proportion to their pages
	\begin{scope}[shift={($(current page.north west)+(0,-9.55cm)$)}]
		\foreach \n/\a/\b in \PackSegments \fill[PackType\n] (\a*21,0) rectangle (\b*21,-0.35);
		\fill[white] (0,0) rectangle (21,-0.05);
		\foreach \n/\a/\b in \PackSegments {\ifdim\a pt>0pt \fill[white] (\a*21-0.04,0) rectangle (\a*21+0.04,-0.35);\fi}
	\end{scope}
	% the spine, split the same way, from under the strip to the footer
	\begin{scope}[shift={($(current page.north west)+(0,-9.9cm)$)}]
		\foreach \n/\a/\b in \PackSegments \fill[PackType\n!60] (0.62,-\a*18.3) rectangle (0.70,-\b*18.3+0.08);
	\end{scope}
	\coverhead
	\coverfooter
\end{tikzpicture}\par
\covertitleblock
% a part: the type colour as a tab, the number, the type, the title, its pages; the whole card is a link
\newcommand{\packcard}[6]{%
	\noindent\hyperlink{bundle-part-#1}{\begin{tikzpicture}[baseline=(t.base)]
		\fill[PackType#1!7, rounded corners=3pt] (0,-0.62) rectangle (17.2,0.62);
		\fill[PackType#1, rounded corners=2pt] (0,-0.62) rectangle (0.24,0.62);
		\node[font=\fontsize{20}{22}\selectfont\bfseries, text=PackType#1] at (0.85,0) {#1};
		\node[anchor=west, font=\scriptsize, text=PackType#1!80!black, inner sep=0pt] at (1.5,0.3) {\spaced{#3}};
		\node[anchor=west, font=\large, text=CoverInk, text width=10.6cm, align=left, inner sep=0pt] (t) at (1.5,-0.2) {#2};
		\node[anchor=east, font=\small\coversans, text=CoverSoft] at (14.5,0) {pages #4--#5};
		\node[anchor=east, font=\small\coversans\bfseries, fill=PackType#1, text=white, rounded corners=4pt,
			inner xsep=5pt, inner ysep=3pt] at (17.0,0) {#6 pages};
	\end{tikzpicture}}\par\vspace{0.22cm}}
\ifdefempty{\PackCards}{}{\infopanel{IN THIS PACK}{\PackCards
	{\small\color{CoverSoft}Tap a document to open it. Each one keeps its own cover, contents and links.\par}}}
\clearpage

% ---------- "In this pack": every part with its outline ----------
% \packmappart{n}{title}{type}{first}{last}: a part's heading, linked to its first page
\newcommand{\packmappart}[5]{%
	\par\vspace{0.35cm}\noindent\hyperlink{bundle-part-#1}{\begin{tikzpicture}[baseline=(t.base)]
		\fill[PackType#1!10] (0,-0.5) rectangle (\linewidth,0.5);
		\fill[PackType#1] (0,-0.5) rectangle (0.18,0.5);
		\node[font=\Large\bfseries, text=PackType#1] at (0.62,0) {#1};
		\node[anchor=west, font=\scriptsize, text=PackType#1!80!black, inner sep=0pt] at (1.15,0.24) {\spaced{#3}};
		\node[anchor=west, font=\bfseries, text=CoverInk, inner sep=0pt] (t) at (1.15,-0.18) {#2};
		\node[anchor=east, font=\small\coversans, text=CoverSoft] at (\linewidth-0.25cm,0) {pages #4--#5};
	\end{tikzpicture}}\par\vspace{0.12cm}}
% \packentry{n}{k}{title}{page}: one line of a part's outline, linked
\newcommand{\packentry}[4]{{\footnotesize\hangindent=1em\hyperlink{bundle-p#1-#2}{\color{CoverInk}#3}\nobreak
	\leaders\hbox to 0.6em{\hss\textcolor{black!30}{.}\hss}\hfill\nobreak
	\hbox to 1.8em{\hfill\hyperlink{bundle-p#1-#2}{\color{PackType#1!80!black}\coversans #4}}\par}}
% \packmapcols{left}{right}: a part's outline in two columns
\newcommand{\packmapcols}[2]{{\setlength{\parskip}{0.06cm}\noindent
	\begin{minipage}[t]{0.475\linewidth}\raggedright#1\end{minipage}\hfill
	\begin{minipage}[t]{0.475\linewidth}\raggedright#2\end{minipage}\par}}
% the back cover lists the parts: \packbackpart{n}{title}{type}{pages}, under a strip split like the front's
\newcommand{\packbackpart}[4]{\noindent\tikz[baseline=-0.6ex]\fill[PackType#1, rounded corners=1pt] (0,-0.22) rectangle (0.34,0.22);\hspace{0.3cm}%
	{\scriptsize\color{PackType#1!80!black}\spaced{#3}}\hspace{0.25cm}{\color{CoverInk}#2}\hfill{\small\coversans\color{CoverSoft}#4 pages}\par\vspace{0.2cm}}
\renewcommand{\CoverBackExtra}{\par\vspace{1.0cm}
	\noindent\begin{tikzpicture}\foreach \n/\a/\b in \PackSegments \fill[PackType\n] (\a*\linewidth,0) rectangle (\b*\linewidth,0.18);
		\foreach \n/\a/\b in \PackSegments {\ifdim\a pt>0pt \fill[white] (\a*\linewidth-0.04cm,0) rectangle (\a*\linewidth+0.04cm,0.18);\fi}\end{tikzpicture}\par
	\vspace{0.4cm}\PackBackParts}
% \packmapmore{n}{text}: what the list leaves out
\newcommand{\packmapmore}[2]{{\small\itshape\color{CoverSoft}#2\par}}
